% source: manti.01.ma, 1ère SM, série fonctions (Drive 0B6pfSJcEVng0ekNJQXNDN3hWMUU)
% Ex. 1 = série française fonctions_s2 (non repris)
\ex{Fonctions}{ekNJ}{Ex. 2}{9}
On considère les fonctions $h(x)=\sqrt{x+2}$ et $g(x)=-\dfrac14x^3$.
\begin{enumerate}
\item Étudier la monotonie de $h$ et de $g$.
\item Tracer les courbes $C_h$ et $C_g$.
\item Montrer que l'équation $x^3+4\sqrt{x+2}=0$ admet une solution dans $\mathbb{R}$.
\item On considère la fonction $f$ définie par $f(x)=\sqrt{2-\dfrac14x^3}$.
\begin{enumerate}
\item Montrer que $D_f=\,]-\infty,2]$.
\item Étudier la monotonie de la fonction $f$.
\end{enumerate}
\end{enumerate}
\fin
\ex{Fonctions}{ekNJ}{Ex. 3}{8}
On pose $h(x)=x^2-2x+3$ et $g(x)=\sqrt{x+3}$.
\begin{enumerate}
\item Étudier la monotonie de $h$ et de $g$.
\item Déterminer les images des intervalles $[-3,-2]$ et $[-2,+\infty[$ par $g$.
\item On considère la fonction $f(x)=x-2\sqrt{x+3}+6$. Vérifier que $f(x)=(h\circ g)(x)$ et étudier la monotonie de $f$.
\end{enumerate}
\fin
\ex{Fonctions}{ekNJ}{Ex. 4}{10}
On considère la fonction $f$ définie par $f(x)=\dfrac{x+1}{\sqrt{x+1}-1}$.
\begin{enumerate}
\item Déterminer l'ensemble de définition de $f$.
\item Montrer que $f$ admet un extremum en $x_0=3$ en précisant sa nature.
\item On pose $g(x)=\dfrac{x^2}{x-1}$.
\begin{enumerate}
\item Montrer que $t_g(x,y)=\dfrac{xy-x-y}{(x-1)(y-1)}$.
\item Étudier la monotonie de $g$ sur les intervalles $]1,2]$ et $[2,+\infty[$ \ ([?] l'original écrit $[0,1[$).
\end{enumerate}
\item Soit $h$ la fonction définie par $h(x)=\sqrt{x+1}$.
\begin{enumerate}
\item Calculer $(g\circ h)(x)$.
\item Étudier la monotonie de la fonction $f$ sur les intervalles $]-1,0[$ et $[3,+\infty[$.
\end{enumerate}
\end{enumerate}
\fin
\ex{Fonctions}{ekNJ}{Ex. 5}{9}
Soit $f$ la fonction définie par $f(x)=\dfrac{2x}{\sqrt{x^2-4}}$.
\begin{enumerate}
\item Déterminer $D_f$ et étudier la parité de $f$.
\item \begin{enumerate}
\item Vérifier que $\big(f(x)\big)^2=4+\dfrac{16}{x^2-4}$.
\item Montrer que $\dfrac{f(x)-f(y)}{x-y}=\dfrac{-16(x+y)}{\left(x^2-4\right)\left(y^2-4\right)}\cdot\dfrac{1}{f(x)+f(y)}$.
\item Étudier la monotonie de $f$ sur $]2,+\infty[$ \ ([?] l'original mentionne aussi $[0,2]$).
\end{enumerate}
\end{enumerate}
\fin
\ex{Fonctions}{ekNJ}{Ex. 6}{8}
On considère la fonction $f$, paire, définie sur $\mathbb{R}$ par :
$f(x)=\sqrt{x-2}$ si $x>2$ \quad et \quad $f(x)=\dfrac{x-2}{x+2}$ si $0\leq x\leq2$.
\begin{enumerate}
\item Calculer $f(-1)$ et $f(-3)$.
\item Dresser le tableau de variation de la fonction $f$ sur $\mathbb{R}$.
\item Tracer la courbe $C_f$.
\item On pose $g(x)=x^2+1$. Déterminer chacune des fonctions $g\circ f$ et $f\circ g$ et dresser leur tableau de variation.
\end{enumerate}
\fin
\ex{Fonctions}{ekNJ}{Ex. 7}{9}
Soit $f$ la fonction définie sur $I=\,]0,1[$ par $f(x)=\left(1+\dfrac1x\right)\left(1+\dfrac{1}{1-x}\right)$.
\begin{enumerate}
\item Montrer que pour tout $x$ de $I$, $1-x$ appartient à $I$ et que $f(1-x)=f(x)$.
\item Vérifier que $f(x)=1-\dfrac{2}{x^2-x}$.
\item Dresser le tableau de variation de $g(x)=x^2-x$ et en déduire la monotonie de $f$.
\item Montrer que si $a$ et $b$ sont deux réels de $\mathbb{R}^+$ avec $a+b=1$, alors $\left(1+\dfrac1a\right)\left(1+\dfrac1b\right)\geq9$.
\end{enumerate}
\fin
